% Chapter 7: Extending Claude — Commands, Skills, Hooks, MCP
\chapter{Extending Claude: Commands, Skills, Hooks, and MCP}
\label{ch:extending}

\begin{itshape}
You keep running the same four commands manually: format, lint, test, commit.
You explain the same deployment procedure every session.
There must be a better way to encode these workflows so Claude executes them
without being told each time.
\end{itshape}

Claude Code ships with four extension mechanisms---commands, skills, hooks,
and MCP servers---forming a hierarchy of increasing autonomy.
Understanding when to use each is the difference between a tool and a workflow.

% -------------------------------------------------------------------
\section{Choosing the Right Mechanism}
\label{sec:extension-strategy}

\marginnote[Tip]{Start with the simplest mechanism.
Promote only when you hit its limits.}

Before documenting each mechanism, here is the strategic question:
\emph{given a workflow problem, which extension do you reach for?}
Getting this wrong costs more than the implementation---it costs
ongoing maintenance on the wrong abstraction.

\subsection*{The Decision Framework}

\begin{decisionbox}[Which Extension Mechanism?]
\textbf{Is the workflow a single prompt that needs a shortcut?}\\
$\rightarrow$ \textbf{Command}. A Markdown file that expands into a prompt.
Zero infrastructure. Examples: \code{/check}, \code{/deploy}, \code{/summarize-pr}.

\textbf{Does it require multi-step orchestration or domain knowledge?}\\
$\rightarrow$ \textbf{Skill}. A prompt-driven playbook that Claude executes
with tool access. Can spawn subagents, read files, make decisions.
Examples: code review workflows, release checklists, research pipelines.

\textbf{Must it happen every time, without exception, regardless of what
Claude is doing?}\\
$\rightarrow$ \textbf{Hook}. Deterministic automation on lifecycle events.
Claude does not decide whether to run it---the harness does.
Examples: auto-format on save, lint on edit, security scan on commit.

\textbf{Does it need to talk to an external service?}\\
$\rightarrow$ \textbf{MCP server}. Gives Claude structured access to databases,
APIs, or tools with schema-defined inputs and outputs.
Examples: Jira, Postgres, Slack, custom internal tools.
\end{decisionbox}

\subsection*{Common Mistakes}

\begin{enumerate}
  \item \textbf{Building a skill when a command suffices.}
    If the workflow is a fixed sequence with no branching logic,
    a command is simpler and easier to maintain. Skills earn their
    complexity when Claude needs to \emph{decide} what to do next.

  \item \textbf{Using a skill when a hook is needed.}
    Skills are advisory---Claude can ignore or adapt them.
    If the behavior must be enforced (security scan, formatting),
    use a hook. Hooks run whether or not Claude remembers to invoke them.

  \item \textbf{Reaching for MCP when a CLI tool exists.}
    \code{gh}, \code{aws}, \code{gcloud}, \code{psql}---Claude already
    knows how to use these. MCP is for services without a good CLI,
    or when you need structured input/output schemas.

  \item \textbf{Encoding team standards in skills instead of hooks.}
    If ``every PR must pass linting'' is a standard, it belongs in a
    \code{PreToolUse} hook on \code{Bash(git commit*)}, not in a
    skill that someone might forget to invoke.
\end{enumerate}

\subsection*{The Promotion Path}

Extensions naturally evolve:

\begin{enumerate}
  \item You type the same prompt three times $\rightarrow$ extract it into a \textbf{command}.
  \item The command needs conditional logic (``if tests fail, do X'')
    $\rightarrow$ promote to a \textbf{skill}.
  \item You realize the skill must run on every commit, not just when invoked
    $\rightarrow$ extract the enforcement part into a \textbf{hook}.
  \item The hook needs data from an external service
    $\rightarrow$ add an \textbf{MCP server} and call it from the hook.
\end{enumerate}

\noindent Each promotion adds capability and maintenance cost.
The simplest mechanism that solves the problem is the best choice.

\begin{keyinsight}
The most common failure mode is not choosing the wrong mechanism---it is
\emph{skipping the command stage entirely}. Most workflows that end up
as complex skills started as prompts that someone typed repeatedly.
If you find yourself building a skill on day one, pause and ask:
would a three-line command file solve this for now?
\end{keyinsight}

% -------------------------------------------------------------------
\section{Custom Slash Commands}
\label{sec:commands}

\marginnote[Official]{Place in \code{\textasciitilde/.claude/commands/} (global)
or \code{.claude/commands/} (project).
Invoked as \code{/command-name}.\sourceurl{https://code.claude.com/docs/en/best-practices}}

\official{Commands are Markdown files that expand into prompts when invoked.\sourceurl{https://code.claude.com/docs/en/skills}}
They accept arguments via \code{\$ARGUMENTS} and can reference files with \code{@path}.

\begin{minted}{markdown}
# .claude/commands/check.md
Run a comprehensive check of the project:
1. Run the test suite: `pytest tests/`
2. Run the linter: `ruff check src/`
3. Run the type checker: `mypy src/`
4. Report results as a summary table
\end{minted}

\practitioner{The best commands are not shortcuts---they encode decisions.}
Instead of a command that runs a test, create a command that decides
\emph{what to do next} based on project state.

\begin{minted}{markdown}
# .claude/commands/next.md
Check the current project state:
1. Read CURRENT_WORK.md for context
2. Run `git status` and `git diff --stat`
3. Check for failing tests
4. Recommend the single highest-priority next action
\end{minted}

\marginnote[Official]{Custom slash commands have been merged into skills.
Both \code{.claude/commands/} and \code{.claude/skills/} work.
Skills take precedence if names conflict.\sourceurl{https://code.claude.com/docs/en/skills}}

% -------------------------------------------------------------------
\section{Skills: Domain Knowledge and Workflows}
\label{sec:skills}

\marginnote[Official]{Each skill lives in its own directory as
\code{SKILL.md}.\sourceurl{https://code.claude.com/docs/en/skills}}

\official{Skills combine domain knowledge with executable workflows.\sourceurl{https://code.claude.com/docs/en/skills}}\index{skills}
Unlike commands (which are prompt templates), skills can:

\begin{itemize}
  \item Be applied automatically when Claude detects a relevant context
  \item Include file globs for automatic triggering
  \item Define multi-step procedures with validation gates
  \item Restrict available tools via \code{allowed-tools} frontmatter
  \item Run shell commands dynamically before sending content to Claude
  \item Spawn a forked subagent context via \code{context: fork}
\end{itemize}

\subsection{File Structure and Precedence}

Each skill is a directory containing a \code{SKILL.md} file:

\begin{minted}{text}
.claude/skills/qa-health/SKILL.md      # Project skill
~/.claude/skills/deploy-check/SKILL.md # Personal (all projects)
\end{minted}

\practitioner{Precedence: enterprise (managed settings)
$>$ personal $>$ project.}
Plugin skills use a namespaced format: \code{plugin-name:skill-name}.

\subsection{Frontmatter Reference}

All frontmatter keys are optional. The \code{name} defaults to the directory name
if omitted.

\begin{fullwidth}
\begin{tabular}{@{}lp{8cm}@{}}
\toprule
\textbf{Key} & \textbf{Purpose} \\
\midrule
\code{name} & Display name and \code{/slash-command}. Lowercase, hyphens, max 64 chars. \\
\code{description} & What the skill does. Claude uses this to decide when to auto-apply. \\
\code{allowed-tools} & Tools permitted without prompting while skill is active. \textbf{Hyphenated.} \\
\code{argument-hint} & Autocomplete hint, e.g.\ \code{[issue-number]}. \\
\code{model} & Override model while skill is active (\code{sonnet}, \code{opus}, \code{haiku}). \\
\code{context} & Set to \code{fork} to run in a forked subagent context. \\
\code{agent} & Which subagent type to use when \code{context: fork} is set. \\
\code{hooks} & Lifecycle hooks scoped to this skill's execution. \\
\code{user-invocable} & Set to \code{false} to hide from \code{/} menu. Default: \code{true}. \\
\code{disable-model-invocation} & Set to \code{true} to prevent Claude from auto-loading. Default: \code{false}. \\
\bottomrule
\end{tabular}
\end{fullwidth}

\subsection{Basic Skill Example}

\begin{minted}{yaml}
---
name: qa-health
description: Run quality dashboard across all volumes
allowed-tools: [Bash, Read, Glob]
---
# QA Health Check
1. Run `python scripts/audits/health_dashboard.py`
2. Parse output for RED/YELLOW/GREEN status
3. Report findings with actionable recommendations
4. If any RED items, create TODO list for fixes
\end{minted}

\subsection{Dynamic Context}

\marginnote[Official]{The \code{!`command`} syntax runs shell commands
at skill load time, injecting live data into the prompt.\sourceurl{https://code.claude.com/docs/en/skills}}

Skills can execute shell commands \emph{before} their content reaches Claude.
This injects live project state into the prompt:

\begin{minted}{yaml}
---
name: review-changes
description: Review recent changes with full context
allowed-tools: [Bash, Read, Grep]
---
# Review Changes

Current branch status:
!`git status --short`

Recent commits:
!`git log --oneline -5`

Review the changes above for logic errors, missing tests,
and style violations per our CLAUDE.md.
\end{minted}

Including ``ultrathink'' anywhere in skill content enables extended thinking
for that skill's execution.

\subsection{Forked Context: Skill as Subagent}

\practitioner{For skills that need isolation---long-running analysis,
exploratory research, or tasks that should not pollute your main context---use
\code{context: fork}.}

\begin{minted}{yaml}
---
name: deep-audit
description: Run thorough code audit in isolated context
context: fork
agent: general-purpose
model: opus
allowed-tools: [Read, Grep, Glob, Bash]
---
# Deep Audit
Analyze the entire codebase for:
1. Security vulnerabilities
2. Performance bottlenecks
3. Dead code
Report findings with file:line references and severity.
\end{minted}

The forked context runs as a subagent---it gets its own context window,
does not consume your main session's context, and returns a summary when done.
This bridges skills and subagents (\cref{ch:agents}).

\subsection{Budget and Permissions}

\marginnote[Official]{Skill descriptions consume $\sim$2\% of context window
(fallback: 16{,}000 characters).
Use \code{/context} to check for excluded skills.\sourceurl{https://code.claude.com/docs/en/skills}}\index{skills!budget (2\%)}

Permission rules can target skills specifically:

\begin{minted}{json}
{
  "permissions": {
    "allow": ["Skill(qa-health)"],
    "deny": ["Skill(deploy *)"]
  }
}
\end{minted}

\code{Skill(name)} matches exactly; \code{Skill(name *)} matches a prefix
with any arguments; bare \code{Skill} denies all skill usage.

\margintip{A skill that saves 5 minutes per use and runs 3 times per week
saves 13 hours per year. Measure value by frequency $\times$ savings.}

% -------------------------------------------------------------------
\section{Hooks: Deterministic Automation}
\label{sec:hooks}

\marginnote[Official]{Configure via \code{/hooks} or in
\code{settings.json}.\sourceurl{https://code.claude.com/docs/en/hooks}}

\convergence{Hooks are the enforcement layer.
Unlike CLAUDE.md instructions (which are advisory),
hooks \textbf{always execute}.}\index{hooks}

\subsection{Available Lifecycle Events (25 Total)}

\marginnote[Official]{Event list and matcher values verified 2026-03-25.
Matchers are specific enum values, not arbitrary
regex.\sourceurl{https://code.claude.com/docs/en/hooks}}\index{hooks!events (24)}

\begin{fullwidth}
\begin{tabular}{@{}lp{4.5cm}p{4cm}@{}}
\toprule
\textbf{Event} & \textbf{Matcher Values} & \textbf{Use Case} \\
\midrule
\code{SessionStart} & \code{startup}, \code{resume}, \code{clear}, \code{compact} & Load context, run health check \\
\code{InstructionsLoaded} & \code{session\_start}, \code{nested\_traversal}, \code{path\_glob\_match}, \code{include}, \code{compact} & Log which instruction files load and when \\
\code{UserPromptSubmit} & (no matcher) & Transform or validate user input \\
\code{PreToolUse} & Tool name (\code{Bash}, \code{Edit}, \code{Write}, \code{Read}, \code{mcp\_\_.*}) & Validate inputs, block operations \\
\code{PermissionRequest} & Tool name & Auto-allow/deny by rule \\
\code{PostToolUse} & Tool name & Auto-lint, log results, trigger follow-up \\
\code{PostToolUseFailure} & Tool name & Error handling, retry logic \\
\code{Notification} & \code{permission\_prompt}, \code{idle\_prompt}, \code{auth\_success}, \code{elicitation\_dialog} & Route to Slack, email, or dashboard \\
\code{Elicitation} & MCP server name & MCP server requests structured user input \\
\code{ElicitationResult} & MCP server name & User responds to an elicitation dialog \\
\code{SubagentStart} & Agent type or name & Initialize subagent-specific config \\
\code{SubagentStop} & Agent type or name & Process subagent results \\
\code{Stop} & (no matcher) & Run cleanup when Claude finishes \\
\code{StopFailure} & \code{rate\_limit}, \code{authentication\_failed}, \code{billing\_error}, \code{server\_error}, \code{max\_output\_tokens}, \code{unknown} & Handle session failures (retry, alert, log) \\
\code{TeammateIdle} & (no matcher) & Notify when agent team member idles \\
\code{TaskCreated} & (no matcher) & React when a background task is spawned \\
\code{TaskCompleted} & (no matcher) & Post-task automation \\
\code{ConfigChange} & \code{user\_settings}, \code{project\_settings}, \code{local\_settings}, \code{policy\_settings}, \code{skills} & React to settings file changes \\
\code{CwdChanged} & (no matcher) & Respond to working directory change \\
\code{FileChanged} & Filename basename & React to external file changes (\code{.envrc}, \code{.env}) \\
\code{WorktreeCreate} & (no matcher) & Custom worktree setup logic \\
\code{WorktreeRemove} & (no matcher) & Custom worktree cleanup logic \\
\code{PreCompact} & \code{manual}, \code{auto} & Save state before context compaction \\
\code{PostCompact} & \code{manual}, \code{auto} & Restore state after compaction completes \\
\code{SessionEnd} & \code{clear}, \code{resume}, \code{logout}, \code{prompt\_input\_exit}, \code{other} & Save state, update logs \\
\bottomrule
\end{tabular}
\end{fullwidth}

\practitioner{The golden rule: if a standard is non-negotiable, make it a hook.
If it is advisory, put it in CLAUDE.md.}\index{hooks!matchers}

\subsection{Handler Types}

\marginnote[Official]{Not all events support all handler types.
8 events support all 4; the remainder are command-only.\sourceurl{https://code.claude.com/docs/en/hooks}}

Hooks support four handler types.\index{hooks!handler types (4)} Command timeout is 10 minutes (increased
from 60 seconds in earlier releases).

\begin{description}
  \item[\code{command}] Shell script. Receives JSON on stdin, returns JSON on stdout.
    The original and most common type.
  \item[\code{http}] POST JSON to a URL (30-second default timeout).
    Best for: webhooks, external service notifications, CI triggers.
  \item[\code{prompt}] Single-turn LLM evaluation returning yes/no.
    Best for: policy checks that need judgment, not just pattern matching.
  \item[\code{agent}] Multi-turn subagent with full tool access.
    Best for: complex validation workflows that read files, run commands,
    or reason across multiple steps.
\end{description}

\begin{minted}{json}
{
  "hooks": {
    "PreToolUse": [{
      "matcher": "Bash",
      "hooks": [
        {
          "type": "command",
          "command": "validate-bash-input.sh"
        },
        {
          "type": "http",
          "url": "https://hooks.internal/audit-log",
          "timeout": 5000
        }
      ]
    }]
  }
}
\end{minted}

\subsection{Notification Hooks}

\official{The \code{Notification} event fires when Claude needs human
attention.\sourceurl{https://code.claude.com/docs/en/common-workflows}}
Use it for desktop alerts during long-running tasks:

\begin{minted}{json}
{
  "hooks": {
    "Notification": [{
      "hooks": [{
        "type": "command",
        "command": "notify-send 'Claude Code' 'Needs attention'",
        "description": "Desktop alert when Claude needs input"
      }]
    }]
  }
}
\end{minted}

\margintip{On macOS, replace \code{notify-send} with
\code{osascript -e 'display notification \"...\"'}.
On Windows WSL, use \code{powershell.exe} with toast notifications.}

\subsection{Debugging Hooks, Skills, and MCP Servers}
\label{subsec:debugging-extensions}

\marginnote[Tip]{Extensions fail silently by default.
Build logging into every hook from the start.}

When an extension does not work and the error is opaque, work through
this checklist:

\begin{description}
  \item[Check if it loaded.]
    Run \code{/hooks} to see merged hook configuration.
    Run \code{/skills} to see available skills.
    Run \code{/mcp} to see connected MCP servers.
    If your extension is missing, the issue is configuration---not logic.

  \item[Check the I/O contract.]
    Command hooks receive JSON on stdin and must return JSON on stdout.
    Exit code 0 = proceed, exit code 2 = block the action.
    Any other exit code or malformed JSON is treated as ``proceed''
    (fail-open). Test your hook standalone:
\begin{minted}{bash}
echo '{"tool":"Bash","input":{"command":"ls"}}' \
  | bash .claude/hooks/my-hook.sh
\end{minted}

  \item[Check permissions.]
    Hook scripts need execute permission (\code{chmod +x}).
    MCP servers may need network access that sandboxing blocks.

  \item[Add logging.]
    Hooks run in the background with no visible output.
    Add \code{>> /tmp/hook-debug.log 2>\&1} to your hook command
    during development. Remove after debugging.

  \item[Check timeouts.]
    Command hooks default to 10 minutes.
    HTTP hooks default to 30 seconds.
    MCP server startup uses \code{MCP\_TIMEOUT} (default varies).
    A slow hook may time out silently.

  \item[Use \code{/debug}.]
    The \code{/debug} skill reads Claude's internal debug log
    and explains what happened. Use it when hooks block
    unexpectedly or MCP connections fail.
\end{description}

\practitioner{The most common hook failure: the script works when
you run it manually but fails in Claude because the PATH is different.
Hook commands run in a minimal shell environment. Use absolute paths
for executables, or source your profile explicitly.}

% -------------------------------------------------------------------
\section{MCP: External Tool Integration}
\label{sec:mcp}

\marginnote[Official]{MCP provides a standardized way to connect Claude to external
tools and data sources.\sourceurl{https://code.claude.com/docs/en/mcp}}

\official{The Model Context Protocol (MCP) connects Claude to
external services, databases, APIs, and custom tools.\sourceurl{https://code.claude.com/docs/en/mcp}}\index{MCP}

\begin{minted}{bash}
# Add an MCP server
claude mcp add github --transport stdio -- gh-mcp

# Scopes
# Local: current machine only (default)
# Project: .mcp.json committed to git (team-wide)
# User: ~/.claude/mcp.json (personal global)
\end{minted}

Key concepts:
\begin{description}
  \item[Servers]\index{MCP!servers} Provide tools, resources, and prompts via a standard protocol.
  \item[Tools]\index{MCP!tools} Functions Claude can call (query a database, create a ticket).
  \item[Resources] Data Claude can reference via \code{@server:resource/path}.
  \item[Dynamic loading] When tools exceed 10\% of context
    (see \cref{sec:context-degradation} for why this matters),
    Claude uses on-demand search to keep context efficient.
\end{description}

\subsection{Practical MCP Walkthrough}

\begin{fullwidth}
\begin{beforeafter}
\before
\begin{minted}{bash}
# Manual: raw API calls to GitHub
curl -H "Authorization: token $GH_TOKEN" \
  https://api.github.com/repos/org/repo/issues
# Parse JSON, extract what you need, paste into Claude
\end{minted}

\after
\begin{minted}{bash}
# One command: Claude can now query GitHub directly
claude mcp add github \
  --transport stdio -- npx @anthropic/mcp-github

# Verify it works
claude mcp list          # Shows: github (stdio)

# Now in any session, Claude can:
#   - List/create issues and PRs
#   - Read repository contents
#   - Search code across repos
\end{minted}
\end{beforeafter}
\end{fullwidth}

\margintip{MCP servers run locally---your credentials never leave your machine.
Claude communicates with the server over stdio, not over the network.}

\official{Permission rules can use wildcards for MCP tool scoping:\sourceurl{https://code.claude.com/docs/en/permissions}}\index{MCP!permissions}

\begin{minted}{json}
{
  "permissions": {
    "allow": ["mcp__github__*"],
    "deny": ["mcp__github__delete_*"]
  }
}
\end{minted}

\margintip{The \code{mcp\_\_server\_\_*} syntax allows all tools from a server.
Combine with specific deny rules for fine-grained control.}

\begin{warnbox}[MCP Security]
Verify server provenance before installation.
Prefer servers from the official registry.
Review permissions each server requires.
Use Managed MCP for enterprise policy enforcement.
\end{warnbox}

\practitioner{Before installing any MCP server, run this 3-step audit:}

\begin{enumerate}
  \item \textbf{Check the README and manifest.}
    What tools does the server expose? Does it request network access,
    filesystem access, or credentials? A GitHub MCP server should not
    need filesystem write access.
  \item \textbf{Review requested permissions.}
    After \code{claude mcp add}, run \code{/permissions} to see what
    the server can do. Deny anything beyond its stated purpose.
    Use \code{mcp\_\_server\_\_*} allows with specific deny rules
    for destructive operations.
  \item \textbf{Test in a sandboxed session.}
    Enable \code{/sandbox} or use a throwaway worktree for the first run.
    Verify the server behaves as documented before granting production access.
\end{enumerate}

\marginnote[Cross-Ref]{For Managed MCP and enterprise policy enforcement,
see \cref{sec:managed-settings}.}

% -------------------------------------------------------------------
\section{Plugins: Community Extensions}
\label{sec:plugins}

\marginnote[Official]{Run \code{/plugin} to browse the marketplace.
Plugins bundle skills, hooks, subagents, and MCP into installable
units.\sourceurl{https://code.claude.com/docs/en/plugins}}

\official{Plugins are pre-packaged extensions from the community and Anthropic.\sourceurl{https://code.claude.com/docs/en/plugins}}\index{plugins}
They bundle multiple extension types into a single installable unit,
avoiding manual configuration.

\begin{description}
  \item[Code intelligence plugins] \official{For typed languages, install a code
    intelligence plugin immediately.\sourceurl{https://code.claude.com/docs/en/discover-plugins}}
    These give Claude precise ``go to definition'' and ``find references''
    capabilities, plus automatic error detection after edits.
  \item[Service integration plugins] Connect to external tools
    (Notion, Figma, monitoring) without writing MCP servers.
\end{description}

\subsubsection{Why Code Intelligence Matters}

Without a code intelligence plugin, Claude uses text-based search
(\code{Grep}, \code{Glob}) to navigate code. This works but produces
false positives and requires reading multiple candidate files.
With a plugin, Claude gets precise symbol navigation:

\begin{fullwidth}
\begin{tabular}{@{}lp{5cm}p{5cm}@{}}
\toprule
\textbf{Operation} & \textbf{Without Plugin} & \textbf{With Plugin} \\
\midrule
Find definition & \code{Grep} for function name, read candidates & Single ``go to definition'' call \\
Find references & \code{Grep} across codebase, filter false positives & Precise ``find references'' \\
Post-edit errors & Run compiler manually or wait for tests & Automatic type error detection \\
\bottomrule
\end{tabular}
\end{fullwidth}

\official{Supported languages: TypeScript, Rust, Go, Java, Python (Pyright),
and more.\sourceurl{https://code.claude.com/docs/en/discover-plugins}}
Install via \code{/plugin} and search for your language.

\marginnote[Cost]{Code intelligence plugins reduce token usage by replacing
broad \code{Grep} searches with precise symbol lookups. Fewer file reads
$=$ smaller context $=$ lower cost per
message.\sourceurl{https://code.claude.com/docs/en/costs}}

\begin{minted}{bash}
# Browse available plugins
/plugin

# Plugins appear in .claude/plugins/ after installation
\end{minted}

\practitioner{Plugins are the fastest path from ``nothing configured''
to ``productive.'' Check the marketplace before building custom skills.}

% -------------------------------------------------------------------
\section{IDE Integrations}
\label{sec:ide-integrations}

\marginnote[Official]{VS Code extension: 2M+ installs.
JetBrains plugin: Beta.\sourceurl{https://code.claude.com/docs/en/vs-code}}

Claude Code runs in three environments: the terminal CLI, the VS Code
extension, and the JetBrains plugin. All share the same underlying
engine, CLAUDE.md files, settings, and hooks---but the interaction
model differs.

\subsection*{VS Code Extension}

\official{The VS Code extension provides a graphical chat panel,
inline diffs with accept/reject controls, and \code{@}-mention file
references.\sourceurl{https://code.claude.com/docs/en/vs-code}}

Install from the Extensions view (\code{Cmd+Shift+X}), search
``Claude Code.'' The extension bundles the CLI---no separate
installation required. Also works with Cursor.

Key capabilities beyond the CLI:

\begin{itemize}
  \item \textbf{Graphical chat panel}: Repositionable in sidebar,
    editor area, or secondary sidebar. Multiple conversations
    via ``Open in New Tab.''
  \item \textbf{Inline diffs}: Side-by-side comparison with
    per-hunk accept/reject controls.
  \item \textbf{\code{@}-mentions}: Type \code{@} for fuzzy file
    matching. \code{@file.ts\#5-10} for line ranges.
    \code{@terminal:name} for terminal output.
  \item \textbf{Plan mode}: Changes displayed as a markdown document
    with inline commenting before approval.
  \item \textbf{Checkpoints}: Hover any message to rewind conversation,
    code, or both.
\end{itemize}

\begin{fullwidth}
{\small
\begin{tabular}{@{}lp{7cm}@{}}
\toprule
\textbf{Shortcut} & \textbf{Action} \\
\midrule
\code{Cmd+Esc} / \code{Ctrl+Esc} & Toggle focus between editor and Claude \\
\code{Cmd+Shift+Esc} & Open Claude in new editor tab \\
\code{Option+K} / \code{Alt+K} & Insert \code{@}-mention from editor \\
\bottomrule
\end{tabular}
}
\end{fullwidth}

\margintip{Set \code{initialPermissionMode} in VS Code settings to
start every conversation in \code{acceptEdits} or \code{plan} mode.}

\subsection*{JetBrains Plugin (Beta)}

\official{The JetBrains plugin supports IntelliJ, PyCharm, WebStorm,
GoLand, Android Studio, and other JetBrains
IDEs.\sourceurl{https://code.claude.com/docs/en/jetbrains}}

Install from the JetBrains Marketplace (search ``Claude Code'').
Unlike VS Code, the JetBrains plugin \emph{requires} the CLI to be
installed separately. The plugin is terminal-based---Claude runs in
the IDE's integrated terminal with IDE features layered on top:

\begin{itemize}
  \item Native JetBrains diff viewer for proposed changes
  \item Current selection and open file shared as context automatically
  \item IDE diagnostics (lint errors, type errors) forwarded to Claude
  \item \code{Cmd+Esc} / \code{Ctrl+Esc} to launch Claude
\end{itemize}

\subsection*{Choosing Your Environment}

\begin{fullwidth}
{\small
\begin{tabular}{@{}lll@{}}
\toprule
\textbf{Feature} & \textbf{VS Code} & \textbf{JetBrains} \\
\midrule
Interface & Graphical chat panel & Terminal + IDE overlays \\
CLI bundled & Yes & No (install separately) \\
Inline diffs & Extension-native & IDE diff viewer \\
\code{@}-mentions & Full (files, terminals, browser) & Files and selection \\
Status & Stable & Beta \\
\bottomrule
\end{tabular}
}
\end{fullwidth}

\convergence{Teams that standardize on one IDE can commit
extension-specific settings to the project. Teams with mixed
IDEs should rely on CLAUDE.md and \code{settings.json}---these
work identically across all three environments.}

% -------------------------------------------------------------------
\section{Sandboxing: Safe Autonomy}
\label{sec:sandboxing}

\marginnote[Official]{Enable with \code{/sandbox}. Restricts filesystem
and network access at the OS level.\sourceurl{https://code.claude.com/docs/en/sandboxing}}

\official{Sandboxing provides OS-level isolation that restricts
filesystem and network access, allowing Claude to work more freely
within defined boundaries.\sourceurl{https://code.claude.com/docs/en/sandboxing}}\index{sandboxing}

With sandboxing enabled, Claude can run commands without per-action
permission prompts because the sandbox prevents access beyond
its boundaries. This is a safer alternative to
\code{--dangerously-skip-permissions}.

% -------------------------------------------------------------------
\section{CLI Tools: Context-Efficient Integration}
\label{sec:cli-tools}

\convergence{CLI tools are the most context-efficient way to interact
with external services.\sourceurl{https://code.claude.com/docs/en/best-practices}}

Install the CLI for services you use regularly:
\begin{itemize}
  \item GitHub: \code{gh} (issues, PRs, comments, releases)
  \item AWS: \code{aws} (S3, Lambda, deployment)
  \item GCP: \code{gcloud} (compute, storage, BigQuery)
  \item Monitoring: \code{sentry-cli}, \code{datadog-ci}
\end{itemize}

Claude knows how to use common CLIs. For unfamiliar ones, try:
\emph{``Use \code{tool --help} to learn about the tool, then use it to solve X.''}

\margintip{CLIs consume far less context than MCP tools or API calls.
Prefer CLIs over MCP for well-known services.}

% -------------------------------------------------------------------
\section{Visual: The Extension Hierarchy}
\label{sec:delegation-hierarchy}

For the decision framework on when to use each mechanism, see
\cref{sec:extension-strategy}.

\begin{figure}[H]
\centering
\begin{tikzpicture}[
  box/.style={draw=WarmBlue, fill=WarmBlue!8, rounded corners=3pt,
    minimum width=3cm, minimum height=0.9cm, font=\small, align=center},
  desc/.style={font=\footnotesize, text=DarkText, align=left},
  arr/.style={->, thick, WarmBlue!60},
  >=Stealth
]
  % Boxes arranged vertically
  \node[box] (cmd) at (0,0) {\textbf{Commands}};
  \node[box] (skill) at (0,1.8) {\textbf{Skills}};
  \node[box, fill=WarmBlue!12] (hook) at (0,3.6) {\textbf{Hooks}};
  \node[box, fill=WarmBlue!16] (mcp) at (0,5.4) {\textbf{MCP Servers}};

  % Descriptions
  \node[desc, anchor=west] at (2.2,0) {Prompt template, inline context\\
    \textit{Use for: simple shortcuts}};
  \node[desc, anchor=west] at (2.2,1.8) {Multi-step workflow, domain knowledge\\
    \textit{Use for: repeated procedures}};
  \node[desc, anchor=west] at (2.2,3.6) {Deterministic lifecycle automation\\
    \textit{Use for: non-negotiable standards}};
  \node[desc, anchor=west] at (2.2,5.4) {External service integration\\
    \textit{Use for: databases, APIs, tools}};

  % Arrows
  \draw[arr] (cmd) -- (skill);
  \draw[arr] (skill) -- (hook);
  \draw[arr] (hook) -- (mcp);

  % Side labels
  \node[font=\footnotesize\itshape, text=WarmBlue!70, rotate=90, anchor=south]
    at (-2.2,2.7) {Increasing capability $\rightarrow$};
  \node[font=\footnotesize\itshape, text=WarmRose!80, rotate=90, anchor=north]
    at (-2.8,2.7) {$\leftarrow$ Increasing setup};
\end{tikzpicture}
\caption{Match extension mechanism to task requirements.
The simplest mechanism that solves the problem is the best choice.}
\label{fig:delegation-hierarchy}
\end{figure}

\begin{keyinsight}
Most developers jump to complex solutions when a simple one would suffice.
A command handles 80\% of workflow needs. Add skills, hooks, and MCP
only when the simpler mechanism is genuinely insufficient.
\cref{ch:agents} explores subagents and parallel delegation patterns
for when even MCP is not enough.
\end{keyinsight}

% -------------------------------------------------------------------
\begin{trybox}
Create a \code{/check} command for your project:
\begin{enumerate}
  \item Create \code{.claude/commands/check.md}.
  \item Include: \code{pytest tests/}, \code{ruff check src/},
    \code{mypy src/}.
  \item Format output as a summary table: Check $|$ Status $|$ Details.
  \item Test it: type \code{/check} in your next session.
\end{enumerate}
Bonus: add a PostToolUse hook matching \code{Edit|Write}: auto-lint written files.
\end{trybox}
